\chapter{Boot and CPU Configuration}

\section{Introduction to Low-Level CPU Control}
In the previous phase, we built a tiny program for a Cortex-M3, ran it in QEMU, and saw it print one line. Several pieces of that program were handed to us as-is: the linker script and the startup code. The primary goal of this phase is to stop trusting those pre-built components, move beyond simply using the toolchain, and take direct control of the CPU[cite: 25]. 

By the end of this phase, developers will understand every instruction the CPU executes between power-on and the first line of \texttt{main()}[cite: 25]. The development process is structured around three core units:
\begin{itemize}
    \item \textbf{Where code and data live:} Understanding the memory layout and directing the linker script to place code and data in the correct physical locations[cite: 25].
    \item \textbf{Taking control of the CPU:} Controlling the boot process through an assembly startup script, configuring faults, and using a debugger[cite: 25].
    \item \textbf{Life without an OS underneath:} Operating in a freestanding environment to comprehend bare-metal limitations, splitting the program into privileged and unprivileged parts, and building a system call (\texttt{svc})[cite: 25].
\end{itemize}

\subsection{Where You Start and Where You End}
All the work in this phase begins on a stripped-down starting skeleton. The original C startup code has been removed. In its place is a new, almost empty assembly file (\texttt{boot/startup.s}). It contains only the two words the CPU absolutely needs at reset (the stack address and the PC), and a \texttt{Reset\_Handler} that does nothing except call \texttt{main()} and loop forever.

This minimal setup exposes a critical physical reality. If we try to run a program that prints a standard string literal and an initialized global character array (which lives in the \texttt{.data} section), the output will be incomplete:

\begin{lstlisting}[language=bash]
make
make run

main: hello from the base
\end{lstlisting}

The second line is missing. Nothing crashes, and no error is printed: the second call simply prints an empty string. Understanding exactly why this happens---and writing the assembly code to fix it---is the primary objective of this chapter. 

Once the boot code, memory initialization, and fault handlers are correctly implemented from scratch, the final program will successfully print:
\begin{lstlisting}
main: privileged hello
main: this text lives in .data
SVC_Handler: privileged code printing on behalf of unprivileged code
BusFault_Handler: a bus fault happened
\end{lstlisting}
\textit{Note: As established previously, QEMU does not exit automatically when the program finishes because the bare-metal image enters an infinite spin loop.}

\section{Memory Technologies in Embedded Systems}
When designing an operating system or bare-metal software, it is crucial to understand the different memory technologies available on the hardware, as this dictates where instructions and variables are placed. 

\subsection{SRAM (Static RAM): The Working Memory}
\begin{itemize}
    \item SRAM sits directly on the CPU's memory bus, meaning every byte has a specific address and can be accessed using standard load/store instructions[cite: 66].
    \item It is extremely fast; on a microcontroller, a single access typically takes about one clock cycle[cite: 66].
    \item It is byte-writable, allowing the CPU to change single bytes at any time[cite: 66].
    \item It is volatile, meaning all data is lost when power is removed, and its contents are entirely undefined (random bits, not zeros) at power-on[cite: 66].
    \item Due to its writable nature, SRAM is the designated location for anything that changes during program execution, such as variables and the stack[cite: 66, 69].
\end{itemize}

\subsection{NOR Flash: Execute-in-Place Memory}
\begin{itemize}
    \item Like SRAM, NOR Flash sits on the memory bus and can be addressed byte by byte, allowing fast random-access reads[cite: 67].
    \item Because it is on the address bus, the CPU can fetch and execute instructions directly from it without needing to copy the program to RAM first[cite: 67]. This feature is known as execute-in-place (XIP)[cite: 67].
    \item It is non-volatile, ensuring the program remains intact after power-off[cite: 67].
    \item Writing to NOR Flash is complex and cannot be done with simple store instructions[cite: 67]. It requires first erasing an entire sector (turning all bits to 1) and then programming specific words (turning bits from 1 to 0), a process that takes microseconds to milliseconds[cite: 67].
    \item Consequently, at runtime, NOR Flash is effectively treated as read-only memory[cite: 67]. It is used to store the code itself and read-only constants (like string literals)[cite: 69].
\end{itemize}

\subsection{NAND Flash and Block Devices}
\begin{itemize}
    \item NAND Flash is the technology used in SSDs, hard disks, and SD cards[cite: 68].
    \item Unlike NOR Flash, it does not sit on the CPU's memory bus and is instead accessed through a separate controller[cite: 68].
    \item Data cannot be addressed byte by byte; it must be accessed in large blocks or pages (e.g., 512 bytes to 4 KB)[cite: 68].
    \item Because it lacks direct memory addressing, the CPU cannot fetch instructions from NAND Flash, nor can it read single variables directly[cite: 68].
    \item To execute a program stored on a NAND device, an OS loader must first copy the blocks into RAM, which is the standard process for desktop computers[cite: 68]. Microcontrollers generally skip this step by running their programs directly from NOR Flash[cite: 68].
\end{itemize}
\subsection{Applying Constraints to C Code}
Given the physical characteristics of these different memory technologies, an operating system must correctly place every element of a C program based on its requirements[cite: 5]. The following table summarizes the hardware constraints:

\begin{table}[h]
\centering
\begin{tabular}{|l|l|l|l|l|}
\hline
\textbf{Technology} & \textbf{On the bus?} & \textbf{Execute from it?} & \textbf{Write at run time?} & \textbf{Keeps data?} \\ \hline
SRAM & yes & yes & yes, bytes & no \\ \hline
NOR flash & yes & yes (XIP) & no (erase/program) & yes \\ \hline
NAND / disk & no (blocks) & no & through a driver & yes \\ \hline
\end{tabular}
\end{table}

\begin{itemize}
    \item \textbf{Code and constants:} Instructions and read-only data (like string literals) should remain in NOR Flash and be executed or read directly from there[cite: 5].
    \item \textbf{Variables (uninitialized):} Standard uninitialized variables that need to be updated at runtime must live in SRAM[cite: 6].
    \item \textbf{Variables (initialized):} Variables defined with an initial value (e.g., \texttt{int x = 5;}) present a unique challenge. The variable lives in RAM, but RAM is random at power-on. Therefore, the initial value \texttt{5} is stored non-volatilely in NOR Flash, and something must explicitly copy it into RAM at boot, before \texttt{main()} runs[cite: 6].
\end{itemize}

\textbf{Crucial Note:} This requirement to manually copy initialized data is the exact reason why the base program from the previous section prints only one line and fails to print the \texttt{.data} string!

\subsection{The Cortex-M3 Memory Map}
To tell the compiler where to place these different elements, a developer must first thoroughly understand the hardware's memory map. On a 32-bit Cortex-M microprocessor, the CPU sees a single 4 GB physical address space[cite: 6]. The regions within this space are fixed by the ARM architecture, while the chip manufacturer decides how much real memory to put in each region[cite: 6]:

\begin{table}[h]
\centering
\begin{tabular}{|l|l|l|}
\hline
\textbf{Address range} & \textbf{Region} & \textbf{On \texttt{mps2-an385}} \\ \hline
\texttt{0x0000\_0000}--\texttt{0x1FFF\_FFFF} & Code & 4 MB at \texttt{0x0}: the program ("FLASH")[cite: 6] \\ \hline
\texttt{0x2000\_0000}--\texttt{0x3FFF\_FFFF} & SRAM & 4 MB at \texttt{0x2000\_0000}: "RAM"[cite: 6] \\ \hline
\texttt{0x4000\_0000}--\texttt{0x5FFF\_FFFF} & Peripherals & UART, timers, \dots[cite: 6] \\ \hline
\texttt{0x6000\_0000}--\texttt{0xDFFF\_FFFF} & External & Not used here[cite: 6] \\ \hline
\texttt{0xE000\_0000}--\texttt{0xE00F\_FFFF} & CPU registers & Control of the CPU itself[cite: 6] \\ \hline
\end{tabular}
\end{table}

Even though we are emulating this hardware in QEMU (where the 4 MB at address 0 is technically just fast RAM pre-loaded with the program image), the emulator enforces these access rules strictly. We treat it as read-only, exactly as we would on a real microcontroller with real flash[cite: 6].

\subsection{The Linker Script: Mapping the Chip}
The \textbf{Linker Script} (e.g., \texttt{scripts/mps2\_m3.ld}) is the file where we describe the hardware's memory map to the toolchain and provide instructions on how to place our code within it[cite: 71].

The script is divided into core sections[cite: 71]:
\begin{itemize}
    \item \texttt{MEMORY}: Defines which physical memory regions exist, their start addresses, and their sizes[cite: 71].
    \item \texttt{SECTIONS}: Instructs the linker on which input sections from the \texttt{.o} files should go into which output regions[cite: 71].
    \item \texttt{ENTRY}: Defines where execution starts[cite: 71].
\end{itemize}

\subsubsection{Defining the Memory Regions}
The \texttt{MEMORY} block dictates the layout of the available hardware. For example[cite: 72]:
\begin{verbatim}
MEMORY {
  FLASH (rx) : ORIGIN = 0x00000000, LENGTH = 4M
  RAM (rwx)  : ORIGIN = 0x20000000, LENGTH = 4M
}
\end{verbatim}
Each region declaration includes an identifier (e.g., \texttt{FLASH}), its starting address (\texttt{ORIGIN}), and its total size (\texttt{LENGTH})[cite: 72]. The linker uses this information to verify that the compiled code fits; if the data assigned to a region exceeds its \texttt{LENGTH}, the link process will fail[cite: 72].

The letters in parentheses specify the access rights for that region[cite: 72]:
\begin{itemize}
    \item \textbf{r:} Read
    \item \textbf{w:} Write
    \item \textbf{x:} Execute
\end{itemize}
For instance, \texttt{(rx)} indicates that the \texttt{FLASH} region can be read and executed, but not written to during normal runtime[cite: 72].

\subsubsection{Defining the Memory Regions (Continued)}
In the \texttt{MEMORY} block, we specify the initial parameters of each physical memory region based on the hardware constraints:
\begin{itemize}
    \item We define the starting address of the \texttt{FLASH} region at \texttt{0x00000000}, which is the predefined start of the code section for Cortex-M3 processors[cite: 72]. 
    \item We define the starting address of the \texttt{RAM} region immediately after, at \texttt{0x20000000}, which is the designated space for static RAM on this architecture[cite: 72].
\end{itemize}

Every address is written in hexadecimal using 32 bits (8 characters). It is imperative that this description perfectly matches the physical hardware specifications. If a developer erroneously marks a region as writable (\texttt{w}) when the underlying technology is strictly read-only (like NOR Flash), the linker will silently generate the ELF file without raising errors. However, the program will inevitably crash at runtime when it attempts an illegal operation[cite: 72].

\subsubsection{Symbol Assignments for the Stack Pointer}
Within the linker script, we can define symbols that the linker will substitute with exact memory addresses during the linking process. A crucial example is calculating the starting address of the stack[cite: 72].

On ARM microprocessors, the stack grows downwards (backwards in memory)[cite: 72, 98]. Therefore, the most efficient approach is to set the initial stack pointer to the very end of the RAM[cite: 72]. In the script, this is done by defining a symbol (e.g., \texttt{\_estack}) and assigning it the value of the RAM's starting address plus its total length[cite: 72]:

\begin{verbatim}
_estack = ORIGIN(RAM) + LENGTH(RAM); 
\end{verbatim}

\subsubsection{Organizing Output Sections}
The core function of the linker script is to declare which sections should exist in the final ELF image and what content should be placed inside them[cite: 73]. 

When compiling a C program composed of multiple source files, the compiler produces multiple object files (\texttt{.o}), each divided into various input sections. The linker script uses rules to aggregate these scattered input sections into unified output sections[cite: 73]. 

For example, to organize the executable code, we define a \texttt{.text} output section[cite: 73]:
\begin{verbatim}
.text : { 
    *(.text*) 
    *(.rodata*) 
    _etext = . 
} > FLASH
\end{verbatim}

\begin{itemize}
    \item \texttt{*(.text*)}: This command tells the linker to examine every input object file (\texttt{*}), find any section whose name starts with \texttt{.text}, and place its contents here[cite: 73]. The wildcard ensures that even if sections are slightly renamed by the compiler, they are correctly included[cite: 73].
    \item \texttt{*(.rodata*)}: This instructs the linker to take all read-only constants from all object files and append them immediately after the code within this same output section[cite: 73].
    \item \texttt{\_etext = .}: The dot (\texttt{.}) represents the location counter (the current memory address)[cite: 73]. This line assigns the current address (which corresponds to the end of the \texttt{.text} and \texttt{.rodata} data) to a symbol named \texttt{\_etext}[cite: 73]. This provides a programmatic way to know exactly where the code section finishes in memory[cite: 73].
    \item \texttt{> FLASH}: Finally, this specifies that the entire \texttt{.text} output section must be placed physically in the region labeled \texttt{FLASH} defined in the \texttt{MEMORY} block[cite: 73].
\end{itemize}


\subsection{Run Address and Load Address: The \texttt{.data} Section}
For read-only data and executable code (the \texttt{.text} section), memory placement is straightforward because they exclusively reside in Flash memory. Code runs exactly where it is stored. 

However, writable global variables that are initialized (e.g., \texttt{int x = 5;}) require a different approach. The \texttt{.data} section needs two distinct addresses:
\begin{lstlisting}[language=make]
_sidata = LOADADDR(.data);
.data :
{
    _sdata = .;
    *(.data*)
    _edata = .;
} > RAM AT > FLASH
\end{lstlisting}

Every output section in the ELF file has two addresses:
\begin{itemize}
    \item \textbf{VMA (Virtual Memory Address / Run Address):} This is where the program \textit{uses} the section when it runs. For variables, this must be in RAM so their values can be updated.
    \item \textbf{LMA (Load Memory Address):} This is where the section's bytes are permanently \textit{stored} in the final binary image, which must be in FLASH.
\end{itemize}
The line \texttt{> RAM AT > FLASH} explicitly dictates this split. 

\subsubsection{The Linker Does Not Copy Data}
The linker script also records the addresses the boot code will need: \texttt{\_sdata} and \texttt{\_edata} define the start and end of \texttt{.data} in RAM, while \texttt{\_sidata} points to the flash address where the stored copy of the initial values begins.

The key point is that \textbf{nothing moves the bytes from the LMA to the VMA by itself}. The linker only records the two addresses, and the CPU certainly does not copy anything at reset. \textbf{Your boot code must manually perform the copy}.

Here is the resulting memory layout mapped by the script:
\begin{lstlisting}
FLASH: [.isr_vector][.text + .rodata][copy of .data] <- _sidata
RAM:   [.data] [.bss] [heap] . . . . . . . [stack v] <- _estack
        ^ _sdata ^ _sbss
\end{lstlisting}

\subsection{BSS, Heap, and Stack Layout}
Other sections of the program must also be defined in the linker script to properly manage the RAM:
\begin{itemize}
    \item \textbf{\texttt{.bss} Section:} This section contains uninitialized global variables. Unlike the \texttt{.data} section, it has no LMA content in Flash; the image only records its size[cite: 75]. It is assigned exclusively to RAM, and the boot code is responsible for filling this memory area with zeros using the \texttt{\_sbss} and \texttt{\_ebss} symbols[cite: 75].
    \item \textbf{Heap and Stack:} If the program utilizes dynamic memory allocation, a heap region must be defined. In standard embedded layouts, the heap grows upwards from the end of the \texttt{.bss} section, while the stack grows downwards from the very end of the RAM (\texttt{\_estack})[cite: 75].
\end{itemize}

\subsubsection{Safety Asserts in the Linker Script}
To ensure system stability, the linker script can include logical conditions using the \texttt{ASSERT} command[cite: 75]. For instance, an \texttt{ASSERT} can be programmed to verify that the predefined limits of the stack (\texttt{\_StackLimit}) and the top of the heap (\texttt{\_heap\_top}) do not overlap[cite: 75]. If this mathematical condition is violated during compilation, the linker will immediately abort the process and throw an error, preventing a catastrophic memory corruption at runtime[cite: 75].

\subsection{The Role of Linker-Defined Symbols}
Symbols defined in the linker script such as \texttt{\_sdata}, \texttt{\_edata}, \texttt{\_sidata}, \texttt{\_sbss}, \texttt{\_ebss}, and \texttt{\_estack} are unusual: they have \textbf{no storage}[cite: 9, 10]. There is no memory cell called \texttt{\_sdata} holding a value. The linker only gives each of them an \textbf{address}, and the address itself \textit{is} the information.

This distinction is critical when using them from C:
\begin{lstlisting}[language=c]
extern uint32_t _sdata, _edata;   /* "there is something at this address" */
uint32_t *start = &_sdata;        /* OK: the address (0x20000000)         */
uint32_t  oops  = _sdata;         /* wrong idea: reads RAM at that address */
\end{lstlisting}

You declare them as \texttt{extern} objects only so that C lets you take their address with \texttt{\&}. Reading \texttt{\_sdata} itself would read whatever happens to be in RAM at \texttt{0x20000000}, which is not what you want. In assembly, the same idea is written \texttt{ldr r0, =\_sdata}, which means \texttt{r0 = \&\_sdata;} in C. The boot code uses these symbols as loop bounds to copy data and zero out the BSS[cite: 9, 10].

\subsubsection{Memory Alignment}
Another critical instruction often found in linker scripts is the \texttt{ALIGN} command. Many RISC processors require instructions and data to be aligned in memory (typically in 4-byte words). If a section ends at an unaligned address, \texttt{ALIGN(4)} inserts padding. Failing to align memory properly will cause the CPU to crash.

\subsection{Reading the Compilation Results}
You now know what the linker script \textit{asks} for. You can check what the linker actually \textit{did} using three tools[cite: 10].

\begin{itemize}
    \item \textbf{The Map File (\texttt{output.map}):} Generated by passing \texttt{-Map=} in the \texttt{Makefile}. It repeats the \texttt{MEMORY} regions and lists every output section, the input sections inside them, and the load address for \texttt{.data}[cite: 10].
    \item \textbf{\texttt{nm} Tool:} Running \texttt{arm-none-eabi-nm -n} lists all symbols sorted by address[cite: 10]. The letter before each name indicates its kind: \texttt{T/t} for text/code, \texttt{R} for read-only data, \texttt{D} for initialized data, \texttt{B} for BSS, and \texttt{A} for absolute symbols[cite: 10].
    \item \textbf{\texttt{objdump} Tool:} Running \texttt{arm-none-eabi-objdump -h} prints one line per output section, showing its size, VMA, and LMA[cite: 10].
\end{itemize}

\subsubsection{Analyzing the Base Image}
Run all three on the base skeleton:
\begin{lstlisting}[language=bash]
make
arm-none-eabi-objdump -h build/scratch-os.axf
arm-none-eabi-nm -n build/scratch-os.axf
\end{lstlisting}

In the \texttt{objdump -h} output, you find \texttt{.isr\_vector} at address \texttt{0x0} with size 8 bytes, and \texttt{.text} starting right after it at \texttt{0x8}. There is no \texttt{.rodata} line because it was folded into \texttt{.text}[cite: 10, 11]. 

The \texttt{.data} line is the most interesting: its VMA is \texttt{0x20000000} (start of RAM), but its LMA is a small flash address (\texttt{0x60} in the base image), right after the code[cite: 10, 11]. In the \texttt{nm} output, \texttt{\_sidata} has exactly that value, \texttt{0x60}. You also find \texttt{\_sdata} at \texttt{0x20000000}, \texttt{\_edata} 32 bytes later (the size of the string literal), \texttt{\_sbss} and \texttt{\_ebss} equal (no zero-initialized variables in the base), and \texttt{\_estack} at \texttt{0x20400000}.

\subsection{Which Section Does a Declaration Go To?}
The compiler decides the \textit{input} section of each thing you declare, and the linker script then decides where that section goes. For declarations at file scope (outside any function), the rules are:

\begin{table}[h]
    \centering
    \begin{tabular}{|l|l|l|}
        \hline
        \textbf{Global declaration (outside any function)} & \textbf{Section} & \textbf{Region} \\ \hline
        a function body & \texttt{.text} & FLASH \\ \hline
        \texttt{const int k = 3;} & \texttt{.rodata} & FLASH \\ \hline
        \texttt{"a string literal"} & \texttt{.rodata} & FLASH \\ \hline
        \texttt{int x = 5;} & \texttt{.data} & RAM, initial value in FLASH \\ \hline
        \texttt{int y;} or \texttt{int z = 0;} & \texttt{.bss} & RAM, zeroed at boot \\ \hline
    \end{tabular}
\end{table}

Note the last row: a variable explicitly initialized to zero also goes to \texttt{.bss}. Because "start at zero" is exactly what \texttt{.bss} guarantees, there is no reason to store the zero in Flash[cite: 11].

\subsubsection{Pointers, Local, and Static Variables}
Pointers deserve a moment of care, because a pointer is itself a variable[cite: 11]. In \texttt{const char *greeting = "hi";}, the characters \texttt{"hi"} form a string literal and go to \texttt{.rodata}[cite: 11]. However, \texttt{greeting}, the variable holding their address, is an initialized writable variable and goes to \texttt{.data}[cite: 11].

The section rules discussed above apply only to declarations made \textbf{outside any function}. A local variable inside a function has no section at all: it lives dynamically on the stack while the function runs, and it is gone as soon as the function returns[cite: 11].

By contrast, a \texttt{static} local variable (e.g., \texttt{static int s = 3;} inside \texttt{main}) behaves like a global variable. It keeps one fixed address for the whole run of the program, exactly like a global---it is only its \textit{name} that stays private to the function. In the \texttt{nm} output, the compiler appends a suffix (like \texttt{s.0}) to distinguish it from other statics with the same name.

\subsubsection{Try It: Analyzing the Output and Little-Endianness}
You can test the section rules by adding these global variables to \texttt{src/main.c}, right after \texttt{data\_message} (before \texttt{main}):
\begin{lstlisting}[language=c]
const int limit = 100;
int counter = 5;
int total;
int zero = 0;
const char *greeting = "hi";
char buffer[64];
\end{lstlisting}

Running \texttt{arm-none-eabi-nm -n build/scratch-os.axf} reveals their placement[cite: 11]. \texttt{counter} and \texttt{greeting} are marked \texttt{D} (initialized data), while \texttt{total}, \texttt{zero}, and \texttt{buffer} are marked \texttt{B} (bss)[cite: 11]. 

Interestingly, in the intermediate object file (\texttt{main.o}), \texttt{limit} is marked \texttt{R} (read-only data), but in the final ELF image, it is marked \texttt{T} (text)[cite: 11, 12]. This happens because the linker script folds \texttt{.rodata} inside the \texttt{.text} output section, and \texttt{nm} classifies a symbol by its final output section[cite: 12]. 

To find the string literal \texttt{"hi"}, which lacks a symbol name, dump the raw contents of the \texttt{.text} output section[cite: 12]:
\begin{lstlisting}[language=bash]
arm-none-eabi-objdump -s -j .text build/scratch-os.axf
\end{lstlisting}
At address \texttt{0x44} you find \texttt{64 00 00 00} (the value 100 for \texttt{limit}), and at \texttt{0x48} you find \texttt{68 69 00} (the characters 'h', 'i', and the null terminator)[cite: 12]. The string \texttt{"hi"} starts exactly at \texttt{0x48} in FLASH.

\paragraph{Pointers in RAM and Little-Endianness}
How does the program find the string? Through \texttt{greeting}, the pointer stored in RAM. Dump the \texttt{.data} section:
\begin{lstlisting}[language=bash]
arm-none-eabi-objdump -s -j .data build/scratch-os.axf
\end{lstlisting}
Following the \texttt{data\_message} string, you will find \texttt{counter} (represented as \texttt{05 00 00 00}) and \texttt{greeting} (represented as \texttt{48 00 00 00}). 

To read these numbers, you must know one crucial hardware fact: the Cortex-M3 is \textbf{little-endian}, meaning a 32-bit number is stored with its lowest byte first. Therefore, the bytes \texttt{48 00 00 00} represent the hexadecimal number \texttt{0x00000048}. This is exactly the FLASH address where \texttt{"hi"} starts. This demonstrates that a pointer is simply a variable whose value is a memory address.

\subsection{Advanced Memory Layout: Custom Sections and Orphans}
You are not limited to the sections the compiler chooses. GCC lets you name the input section of a variable or a function yourself using compiler attributes[cite: 12]:
\begin{lstlisting}[language=c]
int tagged __attribute__((section(".mysection"))) = 42;
\end{lstlisting}

Where that section ends up is up to the linker script, \textbf{if the script names it}[cite: 12]. A section that no rule in the \texttt{SECTIONS} block matches is called an \textbf{orphan}[cite: 12]. The linker does not reject it; it places it by guessing based on its flags (e.g., writable vs. executable) and putting it next to sections that look similar. 

While the guess usually links fine, it places the data outside your control and outside any symbol range you defined. If you add the \texttt{tagged} variable to the base project and search for it:
\begin{lstlisting}[language=bash]
arm-none-eabi-nm -n build/scratch-os.axf | grep -E '_sdata|_edata|tagged'
\end{lstlisting}
The output prints the linker symbols next to your variable:
\begin{lstlisting}
20000000 D _sdata
20000020 D _edata
20000020 D tagged
\end{lstlisting}
Notice that \texttt{tagged} is placed at \texttt{0x20000020}, exactly where \texttt{\_edata} ends. Because it is an orphan, it falls outside the \texttt{\_sdata} to \texttt{\_edata} loop bounds, meaning the boot code will never copy its initial value into RAM!

A new output section \texttt{.mysection} appears. The linker saw a writable, initialized section, so it placed it in RAM \textbf{right after \texttt{.data}}, and gave it its own load address in FLASH. 

Now compare its address with \texttt{\_edata}: \texttt{tagged} sits exactly at \texttt{\_edata}, that is, \textit{after} the end of the range \texttt{\_sdata}...\texttt{\_edata}. A boot loop that copies only \texttt{\_sdata}...\texttt{\_edata} will never copy it, so \texttt{tagged} will not hold \texttt{42} when \texttt{main()} runs. The program links without a single warning and is still wrong. That is the danger of orphans. 

If you declare \texttt{const int} in a section called \texttt{.myconst} instead, the linker sees a read-only section and places it in FLASH, between \texttt{.isr\_vector} and \texttt{.text}.

\subsubsection{Forcing Variables into Specific Sections}
You have the ability to force the compiler to place a variable into a specifically named section directly from your C code. This is done using compiler attributes. For instance:

\begin{verbatim}
int tagged __attribute__((section(".mySeg"))) = 42;
\end{verbatim}

This command defines an integer variable named \texttt{tagged} and forces the linker to store it in a custom section named \texttt{.mySeg}. In the linker script, you can then define exactly where this \texttt{.mySeg} section should be placed physically, giving you absolute control over the memory layout. Even if you do not explicitly define \texttt{.mySeg} in the linker script, the linker will still attempt to manage it as an "orphan" section, though it's best practice to define it explicitly.

\subsection{Overflowing a Region}
The \texttt{LENGTH} values in the \texttt{MEMORY} block are not decoration: the linker enforces them. First, measure how much flash the base image needs:
\begin{lstlisting}[language=bash]
arm-none-eabi-size build/scratch-os.axf
\end{lstlisting}

FLASH has to hold \texttt{text + data}: the code and constants, plus the stored copy of \texttt{.data}'s initial values. For the base skeleton, this is $96 + 32 = 128$ bytes. 

If you edit \texttt{scripts/mps2\_m3.ld} and set FLASH \texttt{LENGTH = 100}, you can test the overflow protection. Rebuild with \texttt{make -B} (because \texttt{make} does not track linker script changes automatically, \texttt{-B} forces everything to be rebuilt)[cite: 12]. The link fails with:
\begin{lstlisting}
region 'FLASH' overflowed by 28 bytes
\end{lstlisting}
and $128 - 100$ is exactly 28[cite: 12]. Setting \texttt{LENGTH = 128} allows the link to succeed because the image fits exactly.

If you restore FLASH and set RAM \texttt{LENGTH = 8K} instead, \texttt{\_estack} moves down to \texttt{0x20002000}. The 4 KB stack reservation starts 4 KB below it, and the 4 KB heap no longer fits between \texttt{.bss} and the stack[cite: 12]. The failure comes from the script's own \texttt{ASSERT} line, which turns a memory corruption that would normally appear at run time into a clear error at build time[cite: 12].

\subsection{Reset: The Hardware Contract}
Once the memory layout is finalized and the ELF file is loaded into memory, the machine can be powered on[cite: 88]. Operating systems on standard desktop computers rely on a BIOS, a loader, and a C runtime to initialize the environment. However, on a bare-metal ARM Cortex-M microprocessor, absolutely none of these services exist[cite: 88].

When the Cortex-M processor comes out of reset, the hardware executes a strictly predefined and immutable sequence[cite: 88]. Before any user code runs, the hardware performs exactly two memory reads at the very beginning of the address space (\texttt{0x00000000})[cite: 88]:
\begin{enumerate}
    \item \textbf{Word 0:} It reads the first 4 bytes at address \texttt{0x0} and copies this value directly into the Stack Pointer (SP) register[cite: 88].
    \item \textbf{Word 1:} It reads the next 4 bytes at address \texttt{0x4} and copies this value into the Program Counter (PC) register, jumping to that address to begin execution[cite: 88].
\end{enumerate}

This behavior is a hardwired contract[cite: 88]. The CPU blindly trusts that the developer has placed the correct initial stack address and the correct starting address of the application code into these first eight bytes of memory[cite: 88]. Everything else—copying the \texttt{.data} section to RAM, zeroing out the \texttt{.bss} section, and setting up the C environment—must be explicitly handled by the developer's boot code starting at the address provided in Word 1[cite: 88].

\subsection{The Vector Table: Controlling the Boot Process}
Those two initial words (SP and PC) are the beginning of the \textbf{vector table}: a table of addresses at a fixed place in memory (address \texttt{0x0} at reset)[cite: 13]. 

The table exists because of \textbf{exceptions}[cite: 13]. An exception is an event that makes the CPU stop what it is doing and run a specific piece of code called its \textbf{handler}[cite: 13]. Reset is one exception; a hardware fault is another[cite: 13]. Each exception has a number, and slot $N$ of the vector table holds the address of the handler for exception $N$[cite: 13]. Slot 0 is special: it is not the address of any code, but the initial value of the SP[cite: 13].

\subsubsection{Implementing the Vector Table in Assembly}
In the base skeleton, the table is written in \texttt{boot/startup.s} inside a section called \texttt{.isr\_vector}, and initially has only two slots[cite: 13]:
\begin{lstlisting}[language=make]
isr_vector:
    .word _estack          /* 0: initial SP  (C: &_estack)      */
    .word Reset_Handler    /* 1: reset       (C: Reset_Handler) */
\end{lstlisting}

\texttt{.word X} simply allocates a 32-bit word and stores the value \texttt{X} inside it[cite: 14, 19]. The label \texttt{isr\_vector} is the address of the first word. Because the linker script places the \texttt{.isr\_vector} section first in FLASH, the table lands exactly at address \texttt{0x0}[cite: 13].

\subsubsection{The Thumb Bit Constraint}
When inspecting the boot sequence in a debugger, the value stored in word 1 is not exactly the raw address of \texttt{Reset\_Handler}, but that address with its lowest bit set to 1. Because instructions always sit at even memory addresses (2-byte or 4-byte aligned), bit 0 of a code address is otherwise unused. ARM uses bit 0 to indicate "this is Thumb code". The Cortex-M executes only the Thumb instruction set, so it strictly requires this bit to be 1[cite: 17, 19].

\subsubsection{\textit{KEEP} and Garbage Collection}
The linker has an optional feature, enabled with \texttt{-Wl,--gc-sections}, called \textbf{garbage collection of sections}. Starting from the \texttt{ENTRY} symbol, the linker follows every reference from one section to another, and \textbf{drops every input section that nothing references}[cite: 14]. This is useful for trimming unused code, but dangerous for data that only the hardware reads: no code references the vector table, so the garbage collector would consider it useless and throw it away[cite: 14].

\texttt{KEEP(...)} in the linker script marks input sections as \textbf{always kept}, whatever the garbage collector thinks[cite: 14]. The rule \texttt{KEEP(*(.isr\_vector))} exists for exactly this reason. 

If you remove \texttt{KEEP} and enable garbage collection, \texttt{objdump -h} will show \texttt{.isr\_vector} with size \texttt{0}. Running the program will fail with a fatal \texttt{Lockup} error. The lockup makes perfect sense: at reset, the CPU still reads words 0 and 1, but those words are now the first instructions of \texttt{Reset\_Handler}, not a stack address and a code address. The CPU jumps into nonsense and stops.

\subsection{Back to the Base: Why Only One Line Prints}
You now have everything you need to explain the base skeleton's behavior from the beginning of the chapter. \texttt{main()} prints two strings:
\begin{enumerate}
    \item The first is a string literal: it lives in \texttt{.rodata}, inside \texttt{.text}, in FLASH, and the CPU reads it directly from there. It prints correctly.
    \item The second is \texttt{data\_message}, an initialized global array. Its VMA is at \texttt{0x20000000} in RAM, but its initial bytes are stored at a flash LMA. 
\end{enumerate}

\texttt{main()} reads the string at its VMA, in RAM. But in the base skeleton, \texttt{Reset\_Handler} contains only \texttt{bl main} and a loop: \textbf{no code moves the bytes from the LMA to the VMA}. The RAM at \texttt{0x20000000} still holds whatever it held at power-on, which in QEMU is all zeros. A string whose first byte is zero is the empty string, so the second call prints nothing! Writing that missing copy is the goal of the next unit.

\section{Taking Control of the CPU}
\subsection{Startup Code: Assembly vs. C}
The previous sections established how to correctly place code and data into memory. The next critical phase is understanding precisely what code must be placed there to initialize the system and take full control of the CPU. The low-level startup code is essential for bridging the gap between hardware power-on and the execution of the main C program.

While this startup sequence can technically be written in C, it is highly instructive—and often more straightforward—to write it directly in Assembly language[cite: 93]. Writing this sequence in C requires complex abstractions and pointer arithmetic to manipulate low-level hardware states[cite: 93]. Assembly, conversely, provides a transparent 1:1 mapping between the instruction written and the physical action taken by the CPU, making the core initialization steps much clearer to understand[cite: 93].

\subsection{What is a Debugger?}
In this unit, you rebuild \texttt{Reset\_Handler} in assembly one step at a time. To do that, you need a debugger to \textit{see} what the CPU does[cite: 14]. A \textbf{debugger} is a program that controls another program while it runs. It can \textbf{stop} the program at any instruction (a \textbf{breakpoint}), \textbf{step} forward one machine instruction at a time, \textbf{inspect} CPU registers and memory, and \textbf{continue} execution[cite: 14].

On a desktop, a crashing program prints an error message. On a microcontroller, there is often no console at all, and a crash simply looks like "nothing happens." A debugger is frequently the \textit{only} way to see what the CPU is really doing[cite: 14]. The standard tool is \textbf{gdb}, the GNU debugger. The ARM version is \texttt{arm-none-eabi-gdb}[cite: 15].

\subsection{A Remote Debugger: Two Programs, Two Terminals}
In this setup, the program being debugged does not run on the same "machine" as GDB[cite: 14]. It runs inside QEMU, which emulates the Cortex-M3 board. GDB talks to QEMU over a network connection (a TCP socket). This requires two separate terminals[cite: 14, 15]:

\begin{verbatim}
 Terminal 1: make debug               Terminal 2: arm-none-eabi-gdb build/scratch-os.axf
 +-------------------------+          +--------------------------+
 | QEMU: CPU + Flash + RAM |  TCP     | gdb: knows your symbols  |
 | halted at reset         |<-------->| (main, Reset_Handler...) |
 | gdb server on port 1234 |  :1234   | from the .axf file       |
 +-------------------------+          +--------------------------+
   program output shows here            you type commands here
\end{verbatim}

\begin{itemize}
    \item \textbf{Terminal 1 (The Target):} You start QEMU with \texttt{make debug}[cite: 15]. The \texttt{debug} target in the \texttt{Makefile} runs QEMU with two crucial options: \texttt{-S}, which freezes the emulated CPU before its very first instruction, and \texttt{-gdb tcp::1234}, which starts a GDB server waiting for a connection on port 1234[cite: 15]. Anything the program prints still appears in this terminal.
    \item \textbf{Terminal 2 (The Host):} You start GDB and pass the \texttt{.axf} ELF file so the debugger has access to all the symbol names[cite: 15]. You type your commands here.
\end{itemize}

\subsubsection{GDB: The GNU Debugger}
The standard tool for this task is \textbf{GDB} (the GNU Debugger)[cite: 94]. While it lacks a graphical user interface and operates entirely via the command line, it is exceptionally powerful and specifically designed for this client-server architecture[cite: 94]. To use it with our ARM target, we run the cross-compiled version (\texttt{arm-none-eabi-gdb}) and load our ELF file (\texttt{scratch-os.axf}) so the debugger has access to all the symbol names (like \texttt{main} or \texttt{Reset\_Handler})[cite: 95]. 

We then connect to the running emulator using the command \texttt{target remote :1234}. From there, GDB allows us to inspect CPU registers (\texttt{info registers}), examine raw memory addresses in hexadecimal (\texttt{x/wx}), step through the program one single machine instruction at a time (\texttt{stepi}), and set breakpoints to pause execution at specific labels (\texttt{break}).

\subsubsection{Setting Up the Debugging Session}
To use the debugger in practice, we must coordinate both ends of the connection:

\textbf{1. Preparing the Target (QEMU):}
Instead of running the emulator normally, we must instruct it to start in debug mode. In the \texttt{Makefile}, there is a specific target named \texttt{debug} which launches QEMU with additional flags:
\begin{verbatim}
-s -gdb tcp::1234
\end{verbatim}
The \texttt{-s} flag is critical: it commands the emulator to freeze the CPU immediately at reset, before the very first instruction is executed[cite: 95]. The \texttt{-gdb tcp::1234} flag opens a GDB server listening on port 1234[cite: 95]. When we type \texttt{make debug}, the terminal will show no program output; the emulator appears stalled because the CPU is powered on but intentionally waiting for a debugger to connect.

\textbf{2. Connecting the Host (GDB):}
In a second terminal, we launch the debugger. It is vital to pass the ELF file as an argument when starting GDB (e.g., \texttt{arm-none-eabi-gdb build/scratch-os.axf})[cite: 95]. GDB does not execute this file; it reads the metadata inside it[cite: 95]. Because we compiled our code with flags that preserve debugging symbols, GDB can map raw hexadecimal memory addresses back to human-readable C variables and function names (like \texttt{main})[cite: 95]. Without the ELF file, the developer would only see raw addresses and would have to manually track the memory locations of every variable.

\textit{Note on Security:} Keeping debugging symbols is essential during development. However, before releasing a final commercial product, these symbols are always stripped from the ELF file. Leaving them in provides malicious actors with a detailed blueprint of how the code operates, making it trivial to exploit.

\subsubsection{Reconstructing the Debug Target}
The \texttt{debug} target is essentially a shorthand, and you can reconstruct it for any image. Take the QEMU command from the \texttt{run} target of the \texttt{Makefile}, and append \texttt{-S -gdb tcp::1234}:
\begin{lstlisting}[language=bash]
qemu-system-arm -machine mps2-an385 -monitor null -semihosting \
    --semihosting-config enable=on,target=native \
    -kernel build/scratch-os.axf -serial stdio -nographic \
    -S -gdb tcp::1234
\end{lstlisting}
This is important for future work: later checkpoints' Makefiles may have a \texttt{run} target but no \texttt{debug} target, and this command is how you will debug them.

\subsubsection{The GDB Commands You Need}
This small set of commands is enough for everything in this phase:
\begin{table}[h]
\centering
\begin{tabular}{|l|l|}
\hline
\textbf{Command} & \textbf{What it does} \\ \hline
\texttt{target remote :1234} & connect to the gdb server (QEMU) on this machine \\ \hline
\texttt{info registers sp pc} & show some registers (\texttt{info registers} alone: all of them) \\ \hline
\texttt{x/2wx 0} & e\textbf{x}amine memory: 2 \textbf{w}ords in he\textbf{x}, at address 0 \\ \hline
\texttt{x/s \&data\_message} & examine memory as a C \textbf{s}tring \\ \hline
\texttt{p/x \$r1} & \textbf{p}rint one register (or a C expression) in hex \\ \hline
\texttt{display/i \$pc} & after every step, show the next \textbf{i}nstruction \\ \hline
\texttt{stepi} (\texttt{si}) & execute exactly one machine instruction (Enter repeats it) \\ \hline
\texttt{break main} (\texttt{b}) & stop when execution reaches \texttt{main} (or any label) \\ \hline
\texttt{continue} (\texttt{c}) & run until the next breakpoint; \texttt{Ctrl-C} interrupts it \\ \hline
\texttt{quit} (\texttt{q}) & leave gdb (QEMU keeps running: stop it with \texttt{Ctrl-C}) \\ \hline
\end{tabular}
\end{table}

\subsubsection{Attaching at Reset}
The first thing to do with the debugger is to catch the CPU \textit{before} its first instruction, and check the reset contract. In the two terminals:
\begin{lstlisting}[language=bash]
make debug                                  # terminal 1: QEMU waits
arm-none-eabi-gdb build/scratch-os.axf      # terminal 2
\end{lstlisting}
And inside GDB:
\begin{lstlisting}
(gdb) target remote :1234
(gdb) info registers sp pc
(gdb) x/2wx 0
\end{lstlisting}

You see \texttt{sp} equal to \texttt{0x20400000}, \texttt{pc} equal to \texttt{0x8 <Reset\_Handler>}, and at address 0 the two words \texttt{0x20400000 0x00000009}. \textbf{No instruction has run yet}, yet SP and PC already have values: the CPU itself copied vector word 0 into SP and vector word 1 into PC. 

Compare the two: word 1 is \texttt{0x9}, but PC is \texttt{0x8}. This is the \textbf{Thumb bit}. The CPU uses bit 0 to check that the target is Thumb code and then clears it; the instruction really starts at \texttt{0x8}. Why \texttt{0x8}? Because the base's vector table has only two words (8 bytes), and the code is placed right after it in FLASH.

\subsection{Cortex-M3 Programmer's Model: The Registers}
To follow the CPU instruction by instruction, you need to know what it holds. The \textbf{programmer's model} is the set of registers your code, and GDB, can see. All of them are 32 bits wide.

\begin{table}[h]
\centering
\begin{tabular}{|l|l|p{8cm}|}
\hline
\textbf{Register} & \textbf{Name} & \textbf{Role} \\ \hline
R0--R12 & general purpose & your values; a C function call passes its first arguments in R0--R3 \\ \hline
R13 & \textbf{SP} & stack pointer: address of the top of the stack \\ \hline
R14 & \textbf{LR} & link register: where to return after a function call \\ \hline
R15 & \textbf{PC} & program counter: address of the next instruction \\ \hline
--- & \textbf{xPSR} & status: flags, Thumb bit, current exception \\ \hline
\end{tabular}
\end{table}

At this level, there are no variables: only these registers and memory. When your C code says \texttt{x = x + 1}, the compiler turns it into "load \texttt{x} from memory into a register, add 1 to the register, store the register back". The stack grows \textbf{downwards}: pushing a value first does \texttt{SP = SP - 4}, then writes the value at the new SP.

\subsubsection{The Status Register (xPSR) and CPU Modes}
\textbf{xPSR} (program status register) packs three critical pieces of information into one register:
\begin{itemize}
    \item \textbf{Flags (Bits 31--28):} The \textbf{N Z C V} flags (negative, zero, carry, overflow), set by comparison instructions and used by conditional jumps.
    \item \textbf{Thumb Bit (Bit 24):} The \textbf{T} state bit, which must always be 1 on Cortex-M. This is where bit 0 of the vector entry ends up.
    \item \textbf{Exception Number (Bits 8--0):} Holds the number of the exception currently being handled, or 0 if none.
\end{itemize}

That last field dictates which of the two \textbf{modes} the CPU is in:
\begin{itemize}
    \item \textbf{Thread mode:} The CPU runs normal code (\texttt{Reset\_Handler}, \texttt{main()}, and the functions they call). The exception number is 0. 
    \item \textbf{Handler mode:} The CPU runs an exception handler, and the field holds that exception's number. 
\end{itemize}
\textit{Note: Code in Thread mode can also be privileged or unprivileged. All code in this specific unit runs privileged; the distinction will be explored in later sections.}
\subsection{ARM Assembly Basics: A Load/Store Architecture}
The boot code must be written in assembly because, before it runs, there is no valid stack and no initialized data, meaning C code cannot yet be trusted. Fortunately, you need only a handful of instructions, and each has a direct C equivalent.

The ARM architecture is strictly a \textbf{load/store machine}[cite: 18]. Memory is touched only by \textbf{load} instructions (\texttt{ldr}, memory to register) and \textbf{store} instructions (\texttt{str}, register to memory)[cite: 18, 19]. Arithmetic and logic operations work only on registers[cite: 18].

\begin{table}[h]
\centering
\begin{tabular}{|l|l|}
\hline
\textbf{Assembly} & \textbf{C equivalent} \\ \hline
\texttt{ldr r3, [r0]} & \texttt{r3 = *r0;} \\ \hline
\texttt{str r3, [r1]} & \texttt{*r1 = r3;} \\ \hline
\texttt{ldr r3, [r0], \#4} & \texttt{r3 = *r0; r0 += 4;} (i.e., \texttt{r3 = *src++;}) \\ \hline
\texttt{str r3, [r1], \#4} & \texttt{*r1 = r3; r1 += 4;} (i.e., \texttt{*dst++ = r3;}) \\ \hline
\texttt{strb r1, [r0, \#7]} & \texttt{((uint8\_t *)r0)[7] = r1;} (one \textbf{b}yte, at offset 7) \\ \hline
\end{tabular}
\end{table}

The square brackets mean "the memory at this address". The form \texttt{[r0], \#4} is called \textbf{post-increment}: the instruction first uses the address in \texttt{r0}, then adds 4 to \texttt{r0}[cite: 22]. Since one word is 4 bytes, a loop using this syntax walks through memory word by word, exactly like \texttt{*dst++} in C.

\subsubsection{Constants, Labels, and Jumps}
Moving data between registers and controlling program flow requires another set of basic instructions:

\begin{table}[h]
\centering
\begin{tabular}{|l|l|}
\hline
\textbf{Assembly} & \textbf{C equivalent} \\ \hline
\texttt{ldr r0, =\_estack} & \texttt{r0 = (uint32\_t)\&\_estack;} (any 32-bit value)[cite: 19] \\ \hline
\texttt{mov sp, r0} & \texttt{sp = r0;} (register to register)[cite: 19] \\ \hline
\texttt{movs r3, \#0} & \texttt{r3 = 0;} (small constants only)[cite: 19] \\ \hline
\texttt{copy\_data:} & a \textbf{label}: a name for this address, like a C \texttt{goto} label \\ \hline
\texttt{b copy\_data} & \texttt{goto copy\_data;} \\ \hline
\texttt{bl main} & \texttt{main();}: \textbf{b}ranch with \textbf{l}ink, saves return address in LR \\ \hline
\end{tabular}
\end{table}

\textbf{The 32-bit Instruction Encoding Constraint:} The \texttt{ldr rX, =value} form needs explanation. An instruction is only 16 or 32 bits long, so it mathematically cannot contain an arbitrary 32-bit constant alongside its opcode[cite: 19]. The assembler solves this by storing the constant in a small table near the code, called the \textbf{literal pool}[cite: 19]. It turns your \texttt{ldr r0, =\_estack} into "load \texttt{r0} from that table". In GDB, this appears as something like \texttt{ldr r0, [pc, \#92]} (a PC-relative address)[cite: 19]. 

The \texttt{.ltorg} directive at the end of the \texttt{Reset\_Handler} tells the assembler where to put this literal pool: typically after the last instruction, where the CPU will never try to execute it as code.

\subsection{Control Flow: Decisions, Bits, and Interrupts}
Branching and loops in assembly rely entirely on the \texttt{xPSR} status register. A decision is a two-step process: compare values to set condition flags, then jump based on those flags[cite: 19].

\begin{table}[h]
\centering
\begin{tabular}{|l|l|}
\hline
\textbf{Assembly} & \textbf{C equivalent} \\ \hline
\texttt{cmp r1, r2} & compute \texttt{r1 - r2}, keep only the flags N Z C V[cite: 20] \\ \hline
\texttt{bhs done} & \texttt{if (r1 >= r2) goto done;} (\textbf{h}igher or \textbf{s}ame, unsigned)[cite: 20] \\ \hline
\texttt{orr r1, r1, \#0x10} & \texttt{r1 |= 0x10;} \\ \hline
\texttt{cpsid i} & stop accepting interrupts (\texttt{\_\_disable\_irq()} in C)[cite: 21] \\ \hline
\texttt{cpsie i} & accept interrupts again (\texttt{\_\_enable\_irq()} in C) \\ \hline
\end{tabular}
\end{table}

With \texttt{cmp} and a conditional branch, you can build any loop. A C \texttt{while} loop becomes "test at the top, jump back at the bottom"[cite: 19, 20]:
\begin{lstlisting}[language=make]
loop:  cmp r1, r2        @ while (r1 < r2) {
       bhs done          @
       ...               @     body
       b   loop          @ }
done:
\end{lstlisting}
The loop continues while \texttt{r1 < r2}; as soon as \texttt{r1 >= r2}, \texttt{bhs} jumps out of the loop[cite: 20]. \textit{(Note: \texttt{@} starts a comment in ARM assembly).}

\subsection{Assembler Directives}
Lines that start with a dot (\texttt{.}) produce no instructions: they steer the assembler itself[cite: 20]. The ones used in \texttt{boot/startup.s} are:

\begin{lstlisting}[language=make]
    .section .isr_vector,"a",%progbits   @ put what follows in section .isr_vector
    .global isr_vector                   @ make the name visible to the linker
isr_vector:
    .word _estack                        @ one 32-bit value
    .word Reset_Handler

    .section .text,"ax",%progbits
    .global Reset_Handler
    .thumb_func                          @ next label is Thumb code (bit 0 = 1)
Reset_Handler:
\end{lstlisting}

This vector table block means, in C equivalent:
\begin{lstlisting}[language=c]
__attribute__((section(".isr_vector")))
const uint32_t isr_vector[] = { (uint32_t)&_estack, (uint32_t)Reset_Handler };
\end{lstlisting}

\begin{itemize}
    \item \texttt{.section} chooses the section, with flags: \texttt{"a"} means the section occupies space in the image, \texttt{"x"} means it is executable.
    \item \texttt{.global} makes a name visible to other files and to the linker; without it, a label acts like a \texttt{static} name in C[cite: 20].
    \item \texttt{.word} allocates exactly one 32-bit word and stores the given value inside it[cite: 20].
    \item \texttt{.thumb\_func} dictates that the next label is Thumb code[cite: 20]. This instructs the tools to store \texttt{Reset\_Handler} with bit 0 set to 1, which the Cortex-M strictly requires[cite: 20].
\end{itemize}
Finally, the file also uses \texttt{.type} and \texttt{.size} directives to tell GDB and \texttt{nm} whether a name is code or data, and how big it is.

\section{Writing the Startup Script (\texttt{Reset\_Handler})}
The \texttt{Reset\_Handler} in \texttt{boot/startup.s} is the very first piece of code executed by the CPU. Its one job is to turn a bare, uninitialized hardware chip into an environment fully capable of safely running C code[cite: 20].

Here is the complete plan, in the order the code does it:
\begin{table}[h]
\centering
\begin{tabular}{|l|l|}
\hline
\textbf{Step} & \textbf{Why} \\ \hline
1. mask interrupts & nothing may interrupt a half-built system \\ \hline
2. set SP & C functions need a stack \\ \hline
3. copy \texttt{.data} FLASH $\rightarrow$ RAM & initialized globals get their values \\ \hline
4. zero \texttt{.bss} & C promises uninitialized globals start at 0 \\ \hline
5. configure the CPU & vector table address, fault handlers, traps, priorities \\ \hline
6. unmask interrupts, \texttt{bl main} & hand over to C \\ \hline
7. loop forever & if \texttt{main} ever returns, there is nowhere to go \\ \hline
\end{tabular}
\end{table}

In the base skeleton, \texttt{Reset\_Handler} has only steps 6 and 7 (and without the unmask). You write everything else. In the rest of this unit, you go through the steps in order, together with the in-class checks that prove each one works.

\subsection{Steps 1 and 2: Mask Interrupts and Set the Stack}
An \textbf{interrupt} is an exception raised by hardware (e.g., a timer expiring or a device receiving data): the CPU drops what it is doing and runs that interrupt's handler[cite: 21]. During setup, nothing is ready to handle one: a handler written in C could read a variable that has not been initialized yet. 

So the very first instruction, \texttt{cpsid i}, tells the CPU to ignore interrupts for now. (No device has been told to raise an interrupt yet, so this is not strictly needed at this moment; writing it makes the rule explicit instead of relying on luck.)

Then, the stack pointer (SP) must be set:
\begin{lstlisting}[language=make]
C:    sp = &_estack;          ASM:  ldr r0, =_estack    @ top of RAM, 0x20400000
                                    mov sp, r0
\end{lstlisting}

C has no way to name SP directly, so the C column is pseudo-code. You might wonder why this is needed, since the hardware already loaded SP from word 0 of the vector table at reset[cite: 21]. The answer is that \texttt{Reset\_Handler} is not always reached through a real reset: a bootloader or a debugger may jump to it directly, with SP holding any random value[cite: 21]. 

Setting SP explicitly makes the code correct in every case, and shows that the stack is nothing more than an address in a register. Notice the two-step pattern: SP is a special register, so the address is first loaded into an ordinary register (\texttt{r0}), and then copied into SP with \texttt{mov}[cite: 21].

\subsection{Step 3: Copy \texttt{.data} from FLASH to RAM}
This is the step that fixes the base skeleton's missing print line[cite: 21]. The linker script gives you three addresses:
\begin{lstlisting}
FLASH:  [ .text | .rodata | initial values of .data ]
                             ^ _sidata
RAM:    [ .data             | .bss | ...        stack ]
          ^ _sdata          ^ _edata                  ^ _estack
\end{lstlisting}

The copy is a simple loop, which you can write first in C and then translate line by line into assembly[cite: 22]:
\begin{lstlisting}[language=make]
C                                     ASM
uint32_t *src = &_sidata;             ldr r0, =_sidata
uint32_t *dst = &_sdata;              ldr r1, =_sdata
                                      ldr r2, =_edata
while (dst < &_edata)             copy_data:  cmp r1, r2
                                              bhs copy_data_done
    *dst++ = *src++;                          ldr r3, [r0], #4
                                              str r3, [r1], #4
                                              b   copy_data
                                  copy_data_done:
\end{lstlisting}

Register \texttt{r0} plays the role of \texttt{src}, \texttt{r1} of \texttt{dst}, and \texttt{r2} holds the end address. The loop is \texttt{copy\_data}, and it exits to \texttt{copy\_data\_done}. It copies whole words, 4 bytes at a time; this is safe because the linker script aligns the start and end of \texttt{.data} to multiples of 4.

\subsubsection{Try It: Bring Data to Life}
In class, you write steps 1-3 into the base skeleton, just before \texttt{bl main}, and run it again:
\begin{lstlisting}[language=bash]
make run
\end{lstlisting}
And the missing line finally appears:
\begin{lstlisting}
main: hello from the base
main: this text lives in .data
\end{lstlisting}
Before the copy, \texttt{data\_message} sat in RAM as all zeros and printed as an empty string; now it successfully holds its initial value[cite: 22].

\subsubsection{Watching the Copy Loop in GDB}
It is worth watching your loop run because it makes the idea of LMA and VMA concrete. Start \texttt{make debug} in terminal 1 and GDB in terminal 2, then:
\begin{lstlisting}
(gdb) target remote :1234
(gdb) display/i $pc
(gdb) stepi                      # press Enter to repeat until "cmp r1, r2"
(gdb) p/x $r0                    # src: a FLASH address
(gdb) p/x $r1                    # dst: 0x20000000, start of RAM
(gdb) x/s &data_message          # "" : nothing copied yet
(gdb) stepi 5                    # one full turn of the loop
(gdb) p/x $r1                    # 0x20000004
(gdb) break copy_data_done
(gdb) continue
(gdb) x/s &data_message          # "main: this text lives in .data\n"
\end{lstlisting}

As you step, \texttt{display/i \$pc} shows each instruction just before it runs. When you reach \texttt{cmp}, \texttt{r1} holds \texttt{0x20000000}, and the string is empty. Five \texttt{stepi} commands execute one full turn of the loop, moving \texttt{r1} to \texttt{0x20000004}. By continuing to the breakpoint at the end, \texttt{r1} climbs until it equals \texttt{r2} (\texttt{0x20000020}), and only then does the string appear in RAM.

\subsubsection{Try It: Breaking the Stack on Purpose}
Seeing a wrong SP go wrong is the best way to believe that SP matters. In \texttt{boot/startup.s}, replace \texttt{\_estack} with a bad address:

\textbf{Test 1: Address Outside of RAM (\texttt{0x30000000})}
If you set SP to \texttt{0x30000000}, nothing prints[cite: 23]. QEMU stops with:
\begin{lstlisting}
qemu: fatal: Lockup: can't escalate 3 to HardFault (current priority -1)
\end{lstlisting}
When \texttt{main} starts, it pushes registers to the stack (below \texttt{0x30000000}), triggering a memory fault[cite: 23]. The CPU tries to run a fault handler to report it, but entering a handler requires writing to the stack, which fails again[cite: 23]. Unable to report the fault, the CPU gives up and enters a \textbf{Lockup} state[cite: 23].

\textbf{Test 2: Unaligned Stack Pointer (\texttt{0x20000001})}
If you force SP to an odd address inside RAM, \texttt{make run} prints nothing[cite: 23]. If you inspect this in GDB (\texttt{info registers sp}), it says \texttt{0x20000000}. The CPU silently drops the two lowest bits of SP because SP must always be a multiple of 4[cite: 23]. The stack now starts at the bottom of RAM and grows \textit{below} it into invalid memory[cite: 23]. The function calls fail silently.

\subsection{Step 4: Zero the \texttt{.bss} Section}
The C standard promises that global variables without an initializer (\texttt{int counter;}) start at 0[cite: 23]. They live in \texttt{.bss}, between \texttt{\_sbss} and \texttt{\_ebss} in RAM, and nothing in flash holds their value: the boot code must actively write the zeros[cite: 23]. 

It is the same loop as the copy, but without a source pointer:
\begin{lstlisting}[language=make]
C                                 ASM
uint32_t *dst = &_sbss;           ldr  r1, =_sbss
                                  ldr  r2, =_ebss
                                  movs r3, #0
while (dst < &_ebss)          zero_bss:  cmp r1, r2
                                         bhs zero_bss_done
    *dst++ = 0;                          str r3, [r1], #4
                                         b   zero_bss
                              zero_bss_done:
\end{lstlisting}

This is \texttt{zero\_bss} in \texttt{boot/startup.s}. In the current program, \texttt{.bss} happens to be empty, so the loop exits immediately, but the first global you add will rely on it. Remember that RAM is \textbf{not} zero at power-on on real hardware; QEMU zeroing it for you is a convenience you must not count on!

\subsection{Exceptions and Faults: One Family}
With steps 1-4, \texttt{Reset\_Handler} does everything that the previous C startup code did. The remaining steps go further: they configure the CPU so that errors are caught and reported.

You have already met the word \textbf{exception}: any hardware event that makes the CPU stop the current code and run a specific handler function in Handler mode[cite: 24]. The Cortex-M3 defines the following family of exceptions:

\begin{table}[h]
    \centering
    \begin{tabular}{|l|l|p{8cm}|}
        \hline
        \textbf{\#} & \textbf{Exception} & \textbf{Raised by} \\ \hline
        0 & \textit{(not an exception)} & this slot holds the initial SP value \\ \hline
        1 & Reset & power-on, reset button \\ \hline
        2 & NMI & non-maskable interrupt: cannot be ignored[cite: 24] \\ \hline
        3 & \textbf{HardFault} & any fault with nowhere else to go[cite: 24] \\ \hline
        4 & \textbf{MemManage} & breaking memory protection rules[cite: 24] \\ \hline
        5 & \textbf{BusFault} & access to an address where nothing answers, or refuses access[cite: 24] \\ \hline
        6 & \textbf{UsageFault} & bad use of an instruction: undefined opcode, Thumb bit 0, divide by 0[cite: 24] \\ \hline
        7--10 & \textit{reserved} & unused: the slot must hold \texttt{0} \\ \hline
        11 & SVCall & the \texttt{svc} instruction \\ \hline
        12 & DebugMonitor & debug events (unused here) \\ \hline
        13 & \textit{reserved} & unused: the slot must hold \texttt{0} \\ \hline
        14 & PendSV & a request raised by software \\ \hline
        15 & SysTick & the CPU's built-in timer \\ \hline
        16+ & IRQs & device interrupts \\ \hline
    \end{tabular}
\end{table}

Apart from the first two slots, all remaining slots in the vector table are associated with these exceptions[cite: 24]. Specific fault handlers (like MemManage, BusFault, and UsageFault) are crucial for both debugging and security[cite: 24].

\subsubsection{Completing the Table with Weak Symbols}
The base table has 2 words, but the CPU expects one word for each exception from 0 to 15 (16 words in total). Every slot needs \textit{some} valid address, even for exceptions you have no handler for yet. The reserved slots (7--10 and 13) must contain \texttt{0}.

To achieve this without writing empty functions for every exception, a best practice is to use a generic template leveraging \textbf{weak symbols} and aliases in assembly:
\begin{lstlisting}[language=make]
    .word HardFault_Handler                    @ slot 3, and so on up to 15
    .weak      HardFault_Handler               @ may be replaced from C
    .thumb_set HardFault_Handler, Default_Handler
Default_Handler:
    b Default_Handler                          @ for (;;) {}
\end{lstlisting}

Which in C (with a GCC extension) translates to:
\begin{lstlisting}[language=c]
void HardFault_Handler(void) __attribute__((weak, alias("Default_Handler")));
\end{lstlisting}

\paragraph{How it works:}
\begin{enumerate}
    \item \texttt{.word HardFault\_Handler} places the address of the label into the table slot.
    \item \texttt{.thumb\_set} tells the assembler: "Make \texttt{HardFault\_Handler} another name for \texttt{Default\_Handler} (same address, marked as Thumb code)."
    \item \texttt{.weak} makes that name a \textbf{weak symbol}: a default that gives way. If a developer writes a normal, non-weak C function named \texttt{void HardFault\_Handler(void)}, the linker silently drops the assembly default and routes the exception directly to the developer's C function.
\end{enumerate}

\texttt{Default\_Handler} itself is an infinite loop: an unexpected exception is a bug, and freezing at a known place is much easier to diagnose with a debugger than running off into random memory.

\paragraph{The Effect on Memory Layout}
Completing the table moves your code. The table grows from 2 words (8 bytes) to 16 words (64 bytes). Because the executable code still comes right after it in FLASH, \texttt{Reset\_Handler} moves from address \texttt{0x8} to \texttt{0x40}. If you attach GDB at reset now, \texttt{info registers pc} shows \texttt{0x40}, and examining address 0 (\texttt{x/2wx 0}) prints \texttt{0x20400000 0x00000041} (the Thumb bit added to \texttt{0x40}).















\subsection{\textcolor{red}{Step 5: Configure the CPU (System Control Block)}}
The remaining configuration steps change settings inside the CPU core itself. You do not need special instructions for most of them: the Cortex-M makes these settings available as 32-bit registers at fixed memory addresses, in the system area of the memory map (\texttt{0xE000E000}--\texttt{0xE000EFFF}). 

The group of registers that controls exceptions is called the \textbf{System Control Block (SCB)}. You change them with ordinary loads and stores:

\begin{lstlisting}[language=make]
C                                   ASM
#define SCB_SHCSR (*(volatile uint32_t *)0xE000ED24)
SCB_SHCSR |= (1u << 16);            ldr r0, =0xE000ED24   @ address
                                    ldr r1, [r0]          @ read
                                    orr r1, r1, #0x10000  @ modify
                                    str r1, [r0]          @ write
\end{lstlisting}

\subsubsection{\textcolor{red}{Read-Modify-Write and the \texttt{volatile} Keyword}}
This pattern is called \textbf{read-modify-write}: you read the current value, set only the bits you care about, and write the result back, so that all the other bits keep their original values.

The \texttt{volatile} keyword in the C version is absolutely critical. It tells the compiler that every memory access must really happen exactly in the order written. The compiler may not skip a read because "the value can't have changed", nor remove a write because "nobody reads it". For hardware registers, both of those optimization assumptions are false.

You can read these registers from GDB like any memory. For example, \texttt{x/wx 0xE000ED00} shows \texttt{0x410fc231}, the CPUID register, which identifies the CPU as an ARM Cortex-M3.

\subsubsection{\textcolor{red}{VTOR: Vector Table Offset Register}}
When an exception happens, the CPU reads the vector table. The \textbf{VTOR} register (at \texttt{0xE000ED08}) holds the memory address of the table. Its value at reset is \texttt{0x00000000}, which is why our initial table must be at address 0. 

However, in real systems, the table is often moved (e.g., a bootloader starting an application higher in flash, or a program copying its table into RAM to modify handlers). Writing VTOR explicitly makes the boot code robust wherever the image is placed:
\begin{lstlisting}[language=make]
C:  *(volatile uint32_t *)0xE000ED08 = (uint32_t)isr_vector;

ASM:  ldr r0, =0xE000ED08
      ldr r1, =isr_vector
      str r1, [r0]
\end{lstlisting}
Because we are setting the entire 32-bit register, no read-modify-write is needed. \textit{(Note: The table address must be aligned to at least 128 bytes).}

\subsubsection{\textcolor{red}{SHCSR: System Handler Control and State Register}}
By default at reset, specific faults (MemManage, BusFault, UsageFault) are disabled, meaning all faults escalate to HardFault. You get a single handler for every kind of error, making it impossible to tell \textit{what} went wrong.

The \textbf{SHCSR} register (at \texttt{0xE000ED24}) enables these specific faults:
\begin{table}[h]
\centering
\begin{tabular}{|l|l|l|}
\hline
\textbf{Bit} & \textbf{Name} & \textbf{When set} \\ \hline
16 & MEMFAULTENA & MemManage faults run \texttt{MemManage\_Handler} \\ \hline
17 & BUSFAULTENA & bus faults run \texttt{BusFault\_Handler} \\ \hline
18 & USGFAULTENA & usage faults run \texttt{UsageFault\_Handler} \\ \hline
\end{tabular}
\end{table}

The boot code must set bits 16--18 (the value \texttt{0x00070000}) using the read-modify-write pattern, because the other bits of SHCSR report status (which handlers are active or pending) and must not be overwritten.




